\documentclass[11pt,a4paper]{article}
\usepackage[T1]{fontenc}
\usepackage{amsmath}
\usepackage{amssymb}
\usepackage{algorithm}
\usepackage{algpseudocode}
\usepackage[margin=2.5cm]{geometry}

\title{Labyrinth Documentation}
\author{}
\date{}

\begin{document}
\maketitle



\section{Labyrinth as a CSP}

\begin{itemize}
\item Goal: Given a parameter size $n$ and a minimum number of walls $nWalls$ (excluding the border), find a solvable labyrinth constellation with an entry and an exit. 
\end{itemize}

Problem Formulation as CSP: 

\[
CSP = \langle \mathcal{X}, \mathcal{D}, \mathcal{C} \rangle
\]

\subsubsection{Variables $\mathcal{X}$}

For a $n \times n$ grid we create $n \times n$ variables, each denoting one cell in the labyrinth grid:

\[
\mathcal{X} = \{X_0, X_1, X_2, \ldots, X_{n \times n - 1}\}
\]

\subsubsection{Domains $\mathcal{D}$}

We denote a wall in the labyrinth as $1$ and an empty cell as $0$. Every variable can either be a wall or an empty cell:

\[
\mathcal{D} = \{D_0, D_1, D_2, \ldots, D_{n \times n - 1}\}
\]
where

\[
D_0 = D_1 = D_2 = \dots = D_{n \times n - 1} = \{0,1\}
\]

\subsubsection{Constraints $\mathcal{C}$}

Each Constraint can be defined as: 

\[
C =
\begin{cases}
Scope(C) & \subseteq \mathcal{X}  \\
Relation(C) & \subseteq \prod_{x \in Sc(C)}^{} D_x
\end{cases}
\]

We differentiate between Tabular Constraints and Predicate Constraints. 
\paragraph{Tabular Constraints}
The Relation, as per definition, is expressed as a set of allowed tuples for the constraint.
\paragraph{Predicate Constraints}
The Relation is expressed as a function. 

Furthermore, constraints can be differentiated based on their \emph{arity}. The arity of a constraint is the number of variables involved in the constraint i.e. the cardinality of its scope (Unary Constraints, Binary Constraints, ..., n-ary Constraints). For CSP-solvers it is desirable to keep the arity small. Any n-ary CSP can be converted into a binary constraint. 

For the labyrinth finder we define the following constraints:

\paragraph{Constraint 1: The top and bottom border of the labyrinth are all walls}
Instead of expressing a single $2n$-ary constraint involving all top and bottom border variables, we can express $2n$ unary constraints. 

First we find all necessary border variables:

\[
\mathcal{X}_{TBborder} = \{X_{B,0}, X_{B,1}, \dots, X_{B,2n}\} \subseteq \mathcal{X}
\]

For each Border variable $X_{B,i}$ we define the following constraint: 

\[
C_{1, i} =
\begin{cases}
Scope(C_{1, i}) = \{X_{B,i}\}  \\
Relation(C_{1, i}) = \{(1)\}
\end{cases}
\]

\begin{itemize}
\item Unary Constraint
\item Tabular Constraint
\end{itemize}

\paragraph{Constraint 2: The left border is all walls except exactly one opening (start).}
Once again we define all left border variables as follows: 

\[
\mathcal{X}_{leftborder} = \{X_{B,0}, X_{B,1}, \dots, X_{B,2n}\} \subseteq \mathcal{X}
\]

Due to the dependence within the variables, we will define a single $2n$-ary constraint.

\[
C_2 =
\begin{cases}
Scope(C_2) = \{X_{B,0}, X_{B,1}, \dots, X_{B,2n}\}  \\
Relation(C_2) = \{(0,1,1, ...), (1, 0, 1, ...), (1, 1, 0, ...), \dots\}
\end{cases}
\]

For the relation, we have exactly $n$ possibilities to distribute one opening among n cells. More precisely this can be expressed as follows: 

\[
Relation(C_2) = \left\{(x_1,\ldots,x_n)\mid \exists j\in\{1,\ldots,n\}:\ x_j=k,\ x_i=0\text{ for all }i\ne j\right\}
\]

\begin{itemize}
\item $2n$-ary Constraint
\item Tabular Constraint
\end{itemize}

\paragraph{Constraint 3: The right border is all wall except exactly one opening (goal).}

We define this constraint equivalently to constraint 2: 

\[
\mathcal{X}_{rightborder} = \{X_{B,0}, X_{B,1}, \dots, X_{B,2n}\} \subseteq \mathcal{X}
\]

\[
C_3 =
\begin{cases}
Scope(C_2) = \{X_{B,0}, X_{B,1}, \dots, X_{B,2n}\}  \\
Relation(C_2) = \{(0,1,1, ...), (1, 0, 1, ...), (1, 1, 0, ...), \dots\}
\end{cases}
\]
or alternatively:
\[
Relation(C_3) = \left\{(x_1,\ldots,x_n)\mid \exists j\in\{1,\ldots,n\}:\ x_j=k,\ x_i=0\text{ for all }i\ne j\right\}
\]

\begin{itemize}
\item $2n$-ary Constraint
\item Tabular Constraint
\end{itemize}

\paragraph{Constraint ?: All open non-border cells have to be reachable from the entrance.}

\paragraph{Constraint ?: Every 2x2 block has to contain at least one open cell. }


\section{Implementation }

\subsection{CSP Solver }

Backtracking search explores partial assignments recursively and returns
when it finds a complete assignment satisfying all constraints.

\subsection{Backtracking search pseudocode}

\begin{algorithm}[htbp]
\caption{Backtracking search for a CSP}
\label{alg:backtracking-search}
\begin{algorithmic}[1]
\Function{Backtracking-Search}{$csp$}
    \State \Return \Call{Backtrack}{$csp, \varnothing$}
\EndFunction
\Statex
\Function{Backtrack}{$csp, assignment$}
    \If{$assignment$ is complete}
        \State \Return $assignment$
    \EndIf
    \State $var \gets \Call{Select-Unassigned-Variable}{csp, assignment}$
    \ForAll{$value$ in \Call{Order-Domain-Values}{$csp, var, assignment$}}
        \If{$value$ is consistent with $assignment$}
            \State add $\{var = value\}$ to $assignment$
            \State $inferences \gets \Call{Inference}{csp, var, assignment}$
            \If{$inferences \neq \mathit{failure}$}
                \State add $inferences$ to $csp$
                \State $result \gets \Call{Backtrack}{csp, assignment}$
                \If{$result \neq \mathit{failure}$}
                    \State \Return $result$
                \EndIf
                \State remove $inferences$ from $csp$
            \EndIf
            \State remove $\{var = value\}$ from $assignment$
        \EndIf
    \EndFor
    \State \Return $\mathit{failure}$
\EndFunction
\end{algorithmic}
\end{algorithm}

Variable selection can use the degree heuristic, domain values can be
ordered by the least-constraining-value heuristic, and inference can use
forward checking, as intended in the project's CSP solver.



\section{Search algorithms to solve a CSP}
\end{document}
